\documentclass[book.tex]{subfiles}

\begin{document}
\chapter{Hilbert Spaces}
\label{chap:hilbert}
\noindent\rule{11cm}{0.4pt} \\
\textbf{Prerequisites:} \autoref{chap:math_basics}, \autoref{chap:calculus}, \autoref{chap:vectors} \\
\textbf{Difficulty Level:} ** \\
\noindent\rule{11cm}{0.4pt}
\vspace{5mm}


Some of the most interest vector spaces in physics and machine learning are infinite dimensional. The words infinite-dimensional can be scary to beginners, so this chapter provides examples of concrete infinite dimensional spaces that are built on familiar mathematical concepts. A Hilbert space is an important special case of infinite dimensional on which there exists a natural inner product.

\section{Spaces of functions}
Let $a, b \in \mathbb{R}$ denote an interval $[a, b]$.
Let $f: [a, b] \to \mathbb{R}$ and $g: [a, b] \to \mathbb{R}$ both be functions. Note then that for $c, d\in \mathbb{R}$, we have that $cf, dg$ are both functions $[a, b] \to \mathbb{R}$. Similarly, $f+g$ is a function $[a, b] \to \mathbb{R}$. Combining these two facts, we have that $cf + dg$ is a function. It follows that the set of such functions forms a vector space. Often times, the space of all functions is too big, so we restrict to continuous functions. We denote the space of all continuous functions from $[a, b]\to \mathbb{R}$ as $C([a, b],\mathbb{R})$.

We can define an inner product on this vector space.

\begin{equation}
\braket{f|g} = \int_a^b f g dx
\end{equation}
The physicists reading along might find this notation evocative. For the rest of our readers, you'll learn more about this "bra-ket" notation later in this book when we cover some quantum mechanics.



\section{Hilbert Spaces}

Let's start off with an abstract definition of a Hilbert space then explain the intuition behind it.


A \textit{Hilbert Space} $\mathcal{H}$ is a complex vector space on which there is an inner product $\langle x | y \rangle$ which associates to each pair of elements $x, y \in \mathcal{H}$ a complex number and which satisfies the following properties

\begin{enumerate}
    \item $\langle y | x \rangle = \overline{\langle x | y \rangle}$: The inner product is conjugate symmetric
    \item $\langle a x_1 + b x_2 | y \rangle = a \langle x_1 | y \rangle + b \langle x_2 | y \rangle$
    \item If $x \neq 0$, then $\langle x | x \rangle > 0$ and if $x = 0$ then $\langle x | x \rangle = 0$
\end{enumerate}

We also require that $\mathcal{H}$ be \textit{complete}, that is limits must exist in the space. Let's consider a more concrete example: The sequence space $\ell^2$ is given by the set of infinite sequences $\mathbf{z} = (z_1, z_2, \dotsc)$ such that
\begin{align}
\sum_{n=1}^\infty |z_n|^2 < \infty
\end{align}
The inner product is defined by
\begin{align}
\langle \mathbf{z} | \mathbf{w} \rangle = \sum_{n=1}^\infty z_n \overline{w_n}
\end{align}

As another example, let $L^2([a, b])$ be the space of functions $f:[a, b] \to \mathbb{R}$ Let the inner product be defined by
\begin{align}
\langle f | g \rangle = \int_a^b f(x)g(x) dx
\end{align}
$L^2([a, b])$ becomes a Hilbert space under this inner product.

%\textcolor{red}{TODO: Add Definition of operators on Hilbert space}

%\textcolor{red}{TODO: Add definition of operator exponential.}

%\textcolor{red}{TODO: Add definitions of eigenvalues/eigenstates}

%\textcolor{red}{TODO: Add more intuition about Hilbert space and size.}

%\textcolor{red}{TODO: Add a discussion of deep learning as choice of function space.}

\section{Exercises}
\begin{enumerate}
    \item Let $\{\left | q\right \rangle \}$ be an orthonormal basis for Hilbert space $\mathcal{H}$. Prove that 
    \begin{equation}
        \sum_i \left | q \right \rangle \left \langle q \right | = I
    \end{equation}
    where $I$ is the identity operator on the Hilbert space $\mathcal{H}$.
    \label{operatoridentity}
    \item Prove that $\ell^2$ is a Hilbert space according to our definition.
    \item Prove that $L^2([a, b])$ is a Hilbert space according to our defnition.
\end{enumerate}


\end{document}